\begin{theorem}
Granica sumy dwóch ciągów zbieżnych jest równa sumie granic tych ciągów, tzn.: jeżeli
$\boldsymbol{\lim\limits_{n\to \infty }a_{n}=a}$ i $\boldsymbol{\lim\limits_{n\to \infty }b_{n}=b}$ to
\[
\boldsymbol{\lim\limits_{n\to \infty }\left (a_{n}+b_{n}\right)=\lim\limits_{n\to \infty }a_{n}+\lim\limits_{n\to \infty }b_{n}=a+b}
\]
\end{theorem}
\begin{proof}[\textbf{Dowód}]
Ponieważ ciągi $\boldsymbol{\left (a_{n} \right)}$ i $\boldsymbol{\left (b_{n} \right)}$ 
są zbieżne to z definicji dla dowolnie małej dodatniej liczby $\boldsymbol{\epsilon}$ 
istnieją takie liczby naturalne $\boldsymbol{N_{1}}$ i $\boldsymbol{N_{2}}$, że dla 
każdej liczby naturalnej $\boldsymbol{n>N_{1}}$ spełniona jest nierówność $\boldsymbol{\left |a_{n}-a \right|<\frac{\epsilon}{2}}$ 
i dla każdej liczby naturalnej $\boldsymbol{n>N_{2}}$ spełniona jest nierówność 
$\boldsymbol{\left | b_{n}-b\right|<\frac{\epsilon}{2}}$. 
Wówczas dla każdego $\boldsymbol{n>N}$, gdzie $\boldsymbol{N}$ to 
największa liczba spośród liczb $\boldsymbol{N_{1}}$ i $\boldsymbol{N_{2}}$ jest:
$\boldsymbol{\left |a_{n}-a \right|+\left |b_{n}-b \right|<\epsilon}$. 
Z własności wartości bezwzględnej wynika, że
\begin{align}
 \boldsymbol{\left |a_{n}-a \right|+\left |b_{n}-b \right|}
 &\boldsymbol{\geqslant \left |a_{n}-a+b_{n}-b \right|} \; \Leftrightarrow \\
\Leftrightarrow\;
&\boldsymbol{\left |a_{n}-a \right|+\left |b_{n}-b \right|\geqslant \left |\left (a_{n}+b_{n} \right)-\left (a+b \right) \right|}, 
\end{align}
a wówczas 
\[
\boldsymbol{\left |\left (a_{n}+b_{n} \right)-\left (a+b \right) \right|\leqslant \left |a_{n}-a \right|+\left |b_{n}-b \right|<\epsilon} 
\]
co daje nam: $\boldsymbol{\left |\left (a_{n}+b_{n} \right)-\left (a+b \right) \right|<\epsilon}$.
\end{proof}